\section{Fundamental Definitions}

\subsection{Latent Activation}
A \textbf{latent activation} (also referred to as a \textit{hidden activation} or \textit{intermediate activation}) denotes the output of a hidden layer \( l \in \{2, \dots, L-1\} \) of a neural network parameterized by \( \hat{f}_{\theta} \), given an input \( x \).
This activation is represented by a vector:
\[
h_l = [a_1, a_2, \dots, a_{n_l}] \in \mathbb{R}^{n_l},
\]
where \( n_l \) denotes the number of neurons in layer \( l \).
We denote by \( \phi_l \) the function mapping an input \( x \) to its latent activation at layer \( l \):
\(
\phi_l : \mathcal{X} \rightarrow \mathbb{R}^{n_l}
\)
The function \( \phi_l \) implicitly contains all transformations applied from the input layer up to layer \( l \).
It can therefore be expressed recursively as the composition of all preceding layers.
At initialization, since the network parameters are randomly initialized, latent activations do not carry any particular meaning.
They merely correspond to random projections of the input data.

\subsection{Latent Representation}
During training, the network parameters are progressively adjusted so that latent activations become informative with respect to the target task.
Consequently, for an input \( x \), the vector \( \phi_l(x) \) evolves in such a way that it captures the relevant characteristics required for the network to produce the expected output.
Once training is completed, these intermediate activations are referred to as \textbf{latent representations} (or \textit{embeddings}).
The term ``representation'' emphasizes that these vectors are no longer simple random projections, but instead encode structured and task-relevant information.
In other words, they constitute a new description of the data, optimized to allow the network \( \hat{f}_{\theta} \) to correctly predict the output \( y \) associated with the input \( x \).

\subsection{Latent Space}
A latent representation only acquires meaning when considered relative to other representations.
Let us consider a dataset:
\[
\{x_1, x_2, \dots, x_m\}.
\]
After training, projecting these samples through the function \( \phi_l \) produces the set:
\[
\{\phi_l(x_1), \phi_l(x_2), \dots, \phi_l(x_m)\}.
\]
This structured set of vectors is called the \textbf{latent space} (also referred to as the \textit{embedding space} or \textit{latent feature space}).
The relative organization of representations within this space reflects how the network distinguishes different inputs and encodes the information necessary for prediction.
It is important to note that the term ``latent space'' should not be understood solely in the mathematical sense of a subspace of \( \mathbb{R}^{n_l} \).
More broadly, it refers to the semantic structure induced by the representations and the relationships between them.
This structure forms the foundation of the interpretation methods studied throughout this chapter.